%%%%%%%%%%%%%%%%%%%%%%%%%%%%%%%%%%%%%%%%%%
\section{Related Work}
\label{sec:related_work}

\subsection{Inertia-Weight Control in Particle Swarm Optimization}
\label{subsec:pso_inertia}

PSO is a population-based stochastic optimizer in which each particle combines inertia with cognitive and social attraction terms \cite{KennedyEberhart1995}. The inertia weight $w$ directly scales the previous velocity; it is therefore a principal mechanism for controlling the exploration--exploitation balance. The original and early variants used a constant coefficient, whereas the widely used linearly decreasing inertia-weight schedule explicitly changes search behavior over the optimization horizon \cite{ShiEberhart1998}. Subsequent work has investigated nonlinear schedules \cite{Chatterjee2006}, feedback from swarm success \cite{Nickabadi2011}, stability-aware adaptation \cite{Taherkhani2015}, and joint adaptation of inertia and acceleration coefficients \cite{Duan2022,Han2024}.

This body of work establishes an important distinction. A time-indexed schedule changes $w$ but does not observe the current search state. State-dependent controllers instead use quantities such as diversity, improvement, distance, or stagnation to react to the population trajectory \cite{Nickabadi2011,Taherkhani2015}. Such feedback is particularly relevant because benchmark landscapes differ substantially in modality and conditioning; a search policy suitable for rapid local refinement need not be suitable for maintaining diversity in a multimodal landscape. The question addressed here is more specific: once a fuzzy state-feedback controller has been selected, how sensitive is its behavior to the semantic resolution of its partitions and to its rule-base structure?

\subsection{Fuzzy State-Feedback Control of PSO Parameters}
\label{subsec:fuzzy_pso}

A fuzzy inference system (FIS) provides an interpretable way to map numerical state descriptors into a control action. In a Mamdani FIS, linguistic antecedents and consequents are combined through a rule base and a defuzzification operator; this architecture remains a canonical model-free fuzzy-control paradigm \cite{Mamdani1974,Lee1990,Nguyen2019Fuzzy}. When used in PSO, the inputs commonly represent population diversity, iteration progress, or an improvement signal, and the output is an adapted value of $w$ or of several PSO coefficients \cite{ShiEberhart2001,Nobile2018,Valdez2017}.

Fuzzy-PSO studies address different control targets and design questions. Valdez et al. compared fuzzy-logic augmentation across PSO variants for mathematical functions using co-evolution \cite{Valdez2017}. Xia et al. designed a multiple-input multiple-output fuzzy controller to tune adaptive learning weights \cite{Xia2023}, while Roy et al. developed benchmark-trained fuzzy feedback policies and reported 127 PSO algorithm variants with up to seven controlled parameters \cite{Roy2024}; the latter article was published in 2025. Salimi et al. applied fuzzy-adaptive PSO to reachability analysis in systems specified through graph transformations \cite{Salimi2025}, a distinct application domain. These contributions demonstrate the breadth of fuzzy parameter control, but they do not change the question addressed here: sensitivity of a fixed static Mamdani inertia architecture to granularity and consequent-pattern choices.

\subsection{Fuzzy Partitions, Linguistic Granularity, and Interpretability}
\label{subsec:granularity}

Linguistic granularity is the number of labels used to discretize a fuzzy universe. It is not merely an implementation detail: it determines which distinctions the controller can express and, for a multi-input rule base, determines the number of candidate rules. A well-formed fuzzy partition is expected to provide coverage, normality, distinguishability, and a semantically ordered set of labels \cite{Klir2005,DuboisPrade2000}. Uniform triangular Ruspini partitions are a common construction because adjacent sets overlap smoothly while their membership grades sum to one over the normalized domain.

The accuracy--interpretability trade-off associated with granularity has long been recognized in fuzzy modeling. Fine partitions can represent smaller input distinctions, but expand the potential rule base; coarser partitions reduce complexity at the cost of local resolution. Gacto et al. review measures for interpretability in linguistic fuzzy rule-based systems \cite{Gacto2011}, while Castiello et al. study variable-granularity interpretable partitions for classified data \cite{Castiello2019}. These works motivate treating granularity as a structural design choice, but they do not make optimization performance a direct measure of interpretability. For the present controller, the number of labels and consequent mapping are separate, coupled choices.

Rule-base design is a second, distinct source of semantic variation. Fuzzy-system design methods can tune membership parameters, induce or prune rules, or modify the partition itself \cite{Ojha2019,Gu2022}. Evolving systems may split rules in response to drift \cite{Lughofer2018} or derive interpretable partitions from observed data \cite{Guillaume2004}. Such adaptive approaches are valuable, but they do not identify whether a particular static combination of granularity and rule logic generalizes over qualitatively different optimization landscapes. Establishing that baseline is necessary before attributing an advantage to more elaborate structural adaptation.

\subsection{From Structural Adaptation to Configuration Selection}
\label{subsec:structural_adaptation}

The broader fuzzy-systems literature distinguishes parameter adaptation from structural adaptation. The former changes membership-function parameters or rule weights; the latter may add, split, merge, or remove rules and partition elements \cite{Gu2022,Saad2020}. Recent research further investigates interpretable evolving partitions \cite{Lughofer2018}, semantic constraints during evolutionary controller design \cite{Megherbi2019}, and variable-granularity knowledge representations \cite{Ahmed2017}. These strands demonstrate that fuzzy semantics need not be treated as immutable.

Nevertheless, the flexibility of an adaptive mechanism can make interpretation difficult when multiple semantic choices change together. A controlled comparison of static configurations provides complementary descriptive evidence about how rankings vary when granularity and rule patterns are changed under a common architecture. Such comparisons do not by themselves identify independent causal effects or confirm a factorial interaction.

\subsection{Research Gap and Positioning of This Study}
\label{subsec:positioning}

Prior fuzzy-PSO studies demonstrate feedback controllers, multi-parameter adaptation, and optimized fuzzy policies, while fuzzy-modeling research treats partition granularity and interpretability as substantive design concerns. The present study is positioned as a controlled structural sensitivity analysis of a static Mamdani fuzzy inertia controller, not as a claim that no earlier factorial study exists. It evaluates four granularities, $n \in \{3,5,7,9\}$, and eight generated consequent patterns, $R_1,\ldots,R_8$, under common inputs, membership-function family, inference method, PSO implementation, and experimental budget. Because granularity changes both partition resolution and rule-table size, the results compare complete configurations rather than identify independent causal effects. The analysis also examines whether performance dispersion among rule-base patterns changes consistently across linguistic granularities, providing descriptive evidence on the sensitivity of complete controller configurations. The study does not validate an adaptive selector.
